% 第10回　提案のブラッシュアップ
\session{10}{2}{12月2日（水）}{提案のブラッシュアップ}{AIによる批判的検証}{yes}

\goals{
  \item 企業が提案を見る観点（課題への適合、実現性、費用対効果、ブランド整合、運用負荷）を挙げられる
  \item AIに批判者の役割を与え、提案の弱点を体系的に洗い出せる
  \item 根拠となるデータの必要性を見極め、AIの数値ではなく一次資料で補強できる
}

\work[100mm]{準備}{提案の要点を1枚にまとめる}
\begin{wstab}{colspec={Q[l,wd=30mm,m] X[l,m]},row{1-Z}={ht=13mm}}
  \hc{企業の課題} & \\ \hc{顧客（ペルソナ）} & \\ \hc{コンセプト} & \\
  \hc{空間} & \\ \hc{体験（動線）} & \\ \hc{期待される効果} & \\ \hc{実現の道筋} & \\
\end{wstab}

\work[80mm]{ワーク1・2}{企業側の視点で批判させ、対応を決める}
\begin{promptbox}[企業側の視点で批判させる]
あなたは自動車ディーラーを運営する会社の経営企画担当者です。次の店舗設計提案について、課題への適合、初期費用と運営費、実現の難しさ、ブランドイメージとの整合、スタッフの運用負荷、安全と法規の観点から、厳しく問題点を指摘してください。指摘は重要な順に並べてください。提案：（要約）
\end{promptbox}
\instr{分類：Ａ＝対応する／Ｂ＝根拠を示して反論する／Ｃ＝割り切る（理由を書く）}
\begin{wstabh}{colspec={X[3,l,m] Q[l,wd=18mm,m] Q[c,wd=10mm,m] X[3,l,m]},row{2-Z}={ht=12mm}}
  AIの指摘（重要な順） & 観点 & 分類 & 具体的な対応・反論の根拠・割り切る理由 \\
  & & & \\ & & & \\ & & & \\ & & & \\ & & & \\ & & & \\
\end{wstabh}

\work[60mm]{根拠}{根拠が示されていない主張を見つける}
\begin{promptbox}[根拠の不足を見つける]
この提案の中で、根拠（データ、事例、調査）が示されていない主張を列挙し、それぞれについて、どんな一次資料で裏づけられそうかを提案してください。数値やデータそのものは生成しないでください。
\end{promptbox}
\begin{wstabh}{colspec={X[3,l,m] X[3,l,m] Q[l,wd=18mm,m]},row{2-Z}={ht=11mm}}
  根拠が不足している主張 & 裏づけに使う一次資料（見つけたら出典） & 担当 \\
  & & \\ & & \\ & & \\ & & \\
\end{wstabh}

\work[60mm]{中間発表}{各グループ5分。受けた指摘と、他グループへのコメント}
\begin{wstabh}{colspec={X[l,m] X[l,m]},row{2-Z}={ht=12mm}}
  自分たちが受けた質問・指摘 & 対応方針 \\
  & \\ & \\ & \\
\end{wstabh}
\answerbox[他グループへのコメント（グループ名と内容）]{28mm}

\begin{nexttime}
  \item 提案の改訂版を作る（指摘への対応を明記）
  \item 発表の構成案のたたき台を作る
\end{nexttime}
\answerbox[改訂のポイント（何をどう変えたか）]{30mm}
\aimemo
